\section{Zastosowanie algorytmu SHADE} \label{sec:shade_application}

Każdy z pięciu wariantów agenta opisanych w sekcji~\ref{sec:modifications} był optymalizowany przy użyciu trzech podejść. Punktem odniesienia jest \textbf{strategia ewolucyjna} w wariancie $(\mu,\lambda) = (10,10)$~\cite{sanchez2020}, w której przystosowanie osobników wyznacza się w schemacie \textbf{koewolucji konkurencyjnej}~\cite{stanley2004} --- osobniki oceniane są nawzajem w obrębie każdej generacji, bez zewnętrznego zbioru referencyjnego. Dwa kolejne podejścia oparto na \textbf{algorytmie SHADE}: wariancie koewolucyjnym oraz wariancie z mechanizmem Hall of Fame. Zestawienie tych trzech optymalizatorów na tych samych wariantach reprezentacji pozwala ocenić, czy adaptacyjna regulacja parametrów SHADE przekłada się na wyższy wskaźnik wygranych względem strategii ewolucyjnej w schemacie koewolucyjnym.

\subsection{SHADE koewolucyjny} \label{sec:shade_coev}

W wariancie koewolucyjnym mutanci i rodzice są oceniani łącznie w każdej generacji. Niech $P_t = \{\mathbf{w}_1, \ldots, \mathbf{w}_{NP}\}$ oznacza populację w generacji $t$, a $U_t = \{\mathbf{u}_1, \ldots, \mathbf{u}_{NP}\}$ --- populację mutantów wygenerowanych przez SHADE. Ocena obejmuje jednocześnie oba zbiory w ramach jednego turnieju w formacie każdy z każdym:
\begin{equation}
    [f_1, \ldots, f_{NP},\; g_1, \ldots, g_{NP}]
    = \textsc{Eval}(P_t \cup U_t)
\end{equation}
gdzie $f_i$ to przystosowanie rodzica $\mathbf{w}_i$, a $g_i$ --- mutanta $\mathbf{u}_i$, obliczone w tej samej puli rozgrywek. Przystosowanie osobnika definiuje się jako łączną liczbę zwycięstw w starciach ze wszystkimi pozostałymi osobnikami w puli:

\begin{equation}
    \phi(\mathbf{w}_i) = \sum_{\mathbf{w}_j \in P_t \cup U_t,\; j \neq i}
        \sum_{d_1 \in \mathcal{D}} \sum_{d_2 \in \mathcal{D}}
        \text{wygr}(\mathbf{w}_i,\; \mathbf{w}_j,\; d_1,\; d_2)
    \label{eq:fitness_coev}
\end{equation}
gdzie $\texttt{NUM\_GAMES}$ to parametr eksperymentalny określający liczbę symulowanych partii między daną parą agentów przy ustalonej kombinacji talii (w niniejszej pracy $\texttt{NUM\_GAMES} = 20$, zestawienie parametrów w sekcji~\ref{sec:params}), a $\text{wygr}(\mathbf{w}_i, \mathbf{w}_j, d_1, d_2)$ oznacza liczbę partii z zakresu $0$ do $\texttt{NUM\_GAMES}$ wygranych przez agenta $\mathbf{w}_i$ w starciu z $\mathbf{w}_j$, gdy gracz pierwszy używa talii $d_1$, a gracz drugi --- $d_2$. Selekcja porównuje każdego mutanta z jego rodzicem na jednolitej skali, eliminując błąd wynikający z losowości środowiska:


\begin{equation}
    \mathbf{w}_i^{(t+1)} =
    \begin{cases}
        \mathbf{u}_i & \text{jeśli } g_i \geq f_i \\
        \mathbf{w}_i & \text{w p.p.}
    \end{cases}
\end{equation}

Niezależnie od wyniku selekcji, wartość przystosowania rodzica jest aktualizowana do wartości $f_i$ obliczonej w bieżącej generacji. Zapobiega to sytuacji, w której fitness rodziców i mutantów byłby obliczany w różnych pulach rozgrywek, co prowadziłoby do nieporównywalnych wartości i błędnych decyzji selekcyjnych.

\subsection{SHADE z Hall of Fame} \label{sec:shade_hof}

W wariancie z Hall of Fame każdy osobnik oceniany jest względem ustalonego zbioru referencyjnego $\mathcal{R}$. Fitness przyjmuje postać analogiczną do \eqref{eq:fitness_coev}, lecz suma przebiega po stałym zbiorze przeciwników:
\begin{equation}
    \phi_{\mathcal{R}}(\mathbf{w}) =
        \sum_{\mathbf{r} \in \mathcal{R}}
        \sum_{d_1 \in \mathcal{D}} \sum_{d_2 \in \mathcal{D}}
        \text{wins}(\mathbf{w},\; \mathbf{r},\; d_1,\; d_2)
    \label{eq:fitness_hof}
\end{equation}
gdzie $\mathbf{w}$ to oceniany osobnik, a $\mathbf{r}$ --- agent ze zbioru referencyjnego $\mathcal{R}$.

Skład zbioru $\mathcal{R}$ dobierano osobno dla każdego wariantu agenta. Dla wariantów korzystających z 21 parametrów (agent bazowy oraz wariant z przeszukiwaniem głębokościowym) zbiór zainicjalizowano 29 agentami:
\begin{itemize}
    \item 20 najlepszych osobników z ostatniej generacji eksperymentu García-Sáncheza i in.~\cite{sanchez2020}, pobranych z repozytorium autorów,
    \item 3 agentów wytrenowanych przez autora niniejszej pracy (200 generacji, $\texttt{POP\_SIZE} = 8$, $\texttt{NUM\_GAMES} = 8$),
    \item 6 agentów o losowo zainicjalizowanych wagach.
\end{itemize}
Dla wariantów z 28 i 63 parametrami przygotowano osobne zbiory referencyjne, złożone z 20 najlepszych osobników uzyskanych w odpowiadającym wariancie eksperymentów koewolucyjnych, przeprowadzonych przed badaniami wariantów SHADE, 3 agentów wytrenowanych na lokalnej maszynie przez autora niniejszej pracy oraz 6 agentów o losowo zainicjalizowanych wagach.

Skład zbioru $\mathcal{R}$ dobrano tak, aby nadać kierunek zbieżności populacji. Rozmiar $|\mathcal{R}| = 29$ nie jest przypadkowy --- w wariancie koewolucyjnym każdy osobnik rywalizuje z pozostałymi $2 \cdot NP - 1 = 29$ osobnikami puli $P_t \cup U_t$, więc zbiór referencyjny o tym samym rozmiarze zapewnia porównywalną liczbę pojedynków i zbliżoną skalę fitness między obiema wersjami. Silni przeciwnicy — wytrenowani wcześniej agenci — zwiększają presję selekcyjną w kierunku skutecznych strategii i ograniczają przedwczesną zbieżność w kierunku optimum lokalnego. Jednocześnie obecność słabszych przeciwników (agenci z losowymi wagami) oraz agentów o umiarkowanej sile daje algorytmowi wczesny sygnał do nauki, a osobniki, które nie potrafią pokonać nawet łatwych przeciwników, są szybko eliminowane. To porządkuje przeszukiwanie, zanim populacja osiągnie poziom wymagający rywalizacji z najsilniejszymi agentami w zbiorze.

Parametr $K_{\text{HOF}} = 5$ oznacza interwał aktualizacji Hall of Fame wyrażony w generacjach --- co piątą generację wykonywana jest podmiana jednego agenta w $\mathcal{R}$. Wartość dobrano jako kompromis między kosztem obliczeniowym a aktualnością przeciwników. Częstsza aktualizacja wymagałaby kosztownej ponownej ewaluacji całej populacji rodziców (sekcja poniżej), a rzadsza utrzymywałaby w $\mathcal{R}$ zbyt długo przestarzałych rywali. Przy typowym przebiegu treningu SHADE z Hall of Fame ($\approx 35$--$43$ generacji przy $\texttt{POP\_SIZE} = 15$) interwał~5 daje około 7--8 podmian w trakcie jednego uruchomienia, co pozwala stopniowo odświeżyć zbiór referencyjny bez nadmiernego narzutu. Zastępowany jest agent, przeciwko któremu bieżąca populacja zdobyła łącznie najwięcej zwycięstw, uznawany jest za najsłabszego i najmniej wymagającego przeciwnika:

\begin{equation}
    r^* = \arg\max_{\mathbf{r} \in \mathcal{R}}\;
        \sum_{i=1}^{NP} \text{wins}(\mathbf{w}_i,\; \mathbf{r},\;
        \cdot,\; \cdot)
\end{equation}
\begin{equation}
    \mathcal{R} \leftarrow
        (\mathcal{R} \setminus \{r^*\}) \cup
        \left\{\arg\max_{\mathbf{w} \in P_t}
        \phi_{\mathcal{R}}(\mathbf{w})\right\}
\end{equation}

Po każdej aktualizacji $\mathcal{R}$ wszyscy rodzice są ponownie ewaluowani względem nowego zbioru, a zaktualizowane wartości fitness zastępują dotychczasowe w wektorze $\phi$, zapewniając spójność skali przed kolejną generacją mutacji. Dodatkowy koszt tej operacji jest przyczyną niższej wartości \texttt{MAX\_EVALUATIONS} w tym wariancie względem wariantu koewolucyjnego.
